\paragraph{KT-SOURCE-CUBIC-EXTRACTION-SCOPE.}
The exact archived source construction does not yet supply
the extra hypotheses (SJ.6)--(SJ.7).
In \path{GRepresents.lean:140--141}, \path{IsLocPt} requires
only that maximal-ideal elements have norm strictly less than
one.  The actual \path{IsPt} predicate at
\path{RhoKApply.lean:76--83} adds integral entries, residual
congruence, the marking, the away-prime condition, and the
local ordinary-component condition.  Its transport theorem
\path{IsPtOf.lean:124--166} preserves exactly those clauses.
\path{TwoPt}, at \path{ConjStable.lean:94--95}, adds failure
of the upper and lower alternatives; it does not specify a
first-order logarithmic support.

The complete normalization proof
\path{NormEx.lean:155--217} gives a congruence-one conjugation
fixing the complex-conjugation matrix and two marked entries.
The finite-coefficient argument
\path{TanDet.lean:110--138,182--216} is qualitative.
The torus point-extension theorem
\path{ExistsTorusTwistEqLocalPointOfWeightedPres.lean:95--137}
and the covering construction \path{ModelUpper.lean:103--250}
also conclude local points, without an entry-depth bound.

There are genuine first-order extension statements:
\path{ConjStable.lean:79--80} extends a specified point
derivation to characteristic-zero dual numbers, and
lines81--85 lift a compatible single-weight base derivation
to the ambient graded model.  Their complete implementations
in \path{ExistsDualNumberPointOfPointDerivationVanishing.lean}
and \path{ExistsBasePointDerivationOfWeightHomogeneous.lean}
do not assert its higher-order integrability.
The local power-series chart at \path{ConjStable.lean:25--37}
is over the fraction field and permits a nonzero clearing
factor; it is not an integral sparse global chart.

For a nonzero first-order coefficient at depth $d$, (SJ.6)
requires the normalized entry valuation
$d/v_\pi(p)<1/3$, whereas locality only gives positivity.
A ramified coefficient extension multiplies both numerator
and denominator by its ramification index.  Torus conjugation
shifts opposite-weight valuations by opposite amounts and
preserves their sum.  Neither operation supplies the missing
estimate or eliminates all intermediate coefficients.
This audit leaves open whether a particular arithmetic point
can satisfy the conditional extraction criterion; it supplies
no such point and no universal point-to-cubic-jet implication.
